\appendix
\chapter{Solutions to Practice Exercises}

This appendix contains step-by-step solutions and explanations for all practice exercises across Weeks 1 to 12.

\section*{Week 1: Introduction to Python, Variables, Data Types \& Operators}

\subsection*{Detailed Solutions}

\textbf{1. Variable Types}
\begin{itemize}
    \item \textbf{(a)} `10 / 2`: \textbf{`float`}. In Python 3, standard division `/` always returns a floating-point number (`5.0`), even if the result is a whole number.
    \item \textbf{(b)} `10 // 2`: \textbf{`int`}. Floor division `//` truncates the fractional part and returns an integer (`5`).
    \item \textbf{(c)} `"100" + "200"`: \textbf{`str`}. Using `+` on two strings performs concatenation, producing `"100200"`.
    \item \textbf{(d)} `(5 > 3) and (2 == 4)`: \textbf{`bool`}. `5 > 3` evaluates to `True`, and `2 == 4` evaluates to `False`. `True and False` evaluates to `False`.
\end{itemize}

\textbf{2. Arithmetic Logic}
\begin{itemize}
    \item \textbf{(a)} `2 + 3 * 4 ** 2`: Exponentiation has highest precedence ($4^2 = 16$). Next is multiplication ($3 \times 16 = 48$). Finally addition ($2 + 48 = \mathbf{50}$).
    \item \textbf{(b)} `25 % 7`: Modulus gives the remainder of $25 \div 7$. $7 \times 3 = 21$, leaving a remainder of $\mathbf{4}$.
    \item \textbf{(c)} `19 // 4`: Floor division of $19 \div 4 = 4.75$, dropping the decimal gives $\mathbf{4}$.
\end{itemize}

\textbf{3. String Slicing}
\begin{itemize}
    \item \textbf{(a)} `text[0:7]` or `text[:7]` returns `"Machine"`.
    \item \textbf{(b)} `text[7:]` returns `"Learning"`.
    \item \textbf{(c)} `text[-8:]` extracts the last 8 characters, which is `"Learning"`.
    \item \textbf{(d)} `text[::-1]` reverses the string to `"gninraeLEnihcaM"`.
\end{itemize}

\textbf{4. Type Casting \& Input}
\begin{lstlisting}
age_str = "22"
future_age = int(age_str) + 5
print(f"In 5 years, you will be {future_age}")
\end{lstlisting}

\textbf{5. Data Science Application: BMI Calculation}
\begin{lstlisting}
weight = 75
height = 1.75
bmi = round(weight / (height ** 2), 2)
print("Patient BMI:", bmi)  # Output: Patient BMI: 24.49
\end{lstlisting}

\section*{Week 2: Control Flow \& Conditionals}

\subsection*{Detailed Solutions}

\textbf{1. Leap Year Logic}
\begin{lstlisting}
year = 2024
if (year % 4 == 0 and year % 100 != 0) or (year % 400 == 0):
    print(f"{year} is a Leap Year")
else:
    print(f"{year} is NOT a Leap Year")
\end{lstlisting}

\textbf{2. Short-Circuit Evaluation}
The output is \textbf{`Safe`}.
Because `x = 0`, the first operand `x != 0` evaluates to `False`. Due to short-circuit evaluation, Python immediately knows the entire `and` expression is `False` without evaluating the second term `(10 / x > 2)`. Thus, division by zero never occurs, preventing a runtime crash.

\textbf{3. Ternary Operator}
\begin{lstlisting}
result = "Pass" if score >= 50 else "Fail"
\end{lstlisting}

\textbf{4. Random Sampling}
\begin{lstlisting}
import random
d1, d2 = random.randint(1, 6), random.randint(1, 6)
print(f"Die 1: {d1}, Die 2: {d2}, Total: {d1 + d2}")
\end{lstlisting}

\textbf{5. Data Science Application: Anomaly Detection}
\begin{lstlisting}
temp = 31.5
if 18 <= temp <= 28:
    status = "Normal"
elif temp < 18:
    status = "Warning: Low Temp"
else:
    status = "Alert: Overheating"

print(status)  # Output: Alert: Overheating
\end{lstlisting}

\section*{Week 3: Loops \& Iteration}

\subsection*{Detailed Solutions}

\textbf{1. Factorial Calculation}
\begin{lstlisting}
N = 6
factorial = 1
curr = N
while curr > 0:
    factorial *= curr
    curr -= 1

print(f"{N}! = {factorial}")  # Output: 6! = 720
\end{lstlisting}

\textbf{2. Pattern Printing}
\begin{lstlisting}
N = 4
for i in range(1, N + 1):
    for j in range(i):
        print("*", end="")
    print()
\end{lstlisting}

\textbf{3. Sum of Evens}
\begin{lstlisting}
even_sum = 0
for i in range(2, 51, 2):
    even_sum += i

print(f"Sum of evens from 1 to 50: {even_sum}")  # Output: 650
\end{lstlisting}

\textbf{4. Early Stopping (`break`)}
\begin{lstlisting}
scores = [12, 45, 67, 89, 92, 40]
for score in scores:
    if score >= 80:
        print(f"Threshold Met: {score}")  # Prints 89
        break
\end{lstlisting}

\textbf{5. Data Science Simulation: Convergence Check}
\begin{lstlisting}
x = 10.0
iterations = 0
while x >= 0.01:
    x *= 0.5
    iterations += 1

print(f"Converged in {iterations} iterations. Final x: {x}")
# Output: Converged in 10 iterations. Final x: 0.009765625
\end{lstlisting}

\section*{Week 4: Lists \& Basic Matrix Algorithms}

\subsection*{Detailed Solutions}

\textbf{1. List Operations}
\begin{lstlisting}
data = [15, 8, 42, 4, 23]
data.append(16)   # [15, 8, 42, 4, 23, 16]
data.sort()       # [4, 8, 15, 16, 23, 42]
data.reverse()    # [42, 23, 16, 15, 8, 4]
print(data)
\end{lstlisting}

\textbf{2. Vector Euclidean Distance}
\begin{lstlisting}
import math
p = [1, 2]
q = [4, 6]

dist = math.sqrt((p[0] - q[0])**2 + (p[1] - q[1])**2)
print("Euclidean Distance:", dist)  # Output: 5.0
\end{lstlisting}

\textbf{3. Matrix Transpose}
\begin{lstlisting}
A = [[1, 2, 3], 
     [4, 5, 6]]

AT = []
for j in range(len(A[0])):  # Columns of A become rows of AT
    row = []
    for i in range(len(A)):
        row.append(A[i][j])
    AT.append(row)

print("A Transposed:", AT)  # Output: [[1, 4], [2, 5], [3, 6]]
\end{lstlisting}

\textbf{4. Manual Selection Sort}
\begin{lstlisting}
arr = [64, 25, 12, 22, 11]
n = len(arr)

for i in range(n):
    min_idx = i
    for j in range(i + 1, n):
        if arr[j] < arr[min_idx]:
            min_idx = j
    # Swap minimum found with first unsorted element
    arr[i], arr[min_idx] = arr[min_idx], arr[i]

print("Sorted Array:", arr)  # Output: [11, 12, 22, 25, 64]
\end{lstlisting}

\textbf{5. Data Science Application: Min-Max Normalization}
\begin{lstlisting}
features = [10, 20, 30, 40, 50]
x_min = min(features)
x_max = max(features)

scaled = []
for x in features:
    x_scaled = (x - x_min) / (x_max - x_min)
    scaled.append(round(x_scaled, 2))

print("Normalized Features:", scaled)
# Output: [0.0, 0.25, 0.5, 0.75, 1.0]
\end{lstlisting}

\section*{Week 5: Tuples, Dictionaries \& Sets}

\subsection*{Detailed Solutions}

\textbf{1. Tuple Operations}
\begin{lstlisting}
student = ("Alice", 21, "Data Science", 3.8)
name, age, major, gpa = student

# Trying to modify student[1] = 22 raises:
# TypeError: 'tuple' object does not support item assignment
\end{lstlisting}

\textbf{2. Dictionary Lookup \& Updating}
\begin{lstlisting}
stock_prices = {"AAPL": 175.5, "GOOGL": 140.2, "MSFT": 380.0}

stock_prices["AAPL"] = 180.0
stock_prices["AMZN"] = 150.0

nvda_price = stock_prices.get("NVDA", 0.0)
print("NVDA Price:", nvda_price)  # Output: 0.0
\end{lstlisting}

\textbf{3. Set Membership \& Deduplication}
\begin{lstlisting}
ids = [101, 102, 101, 105, 102, 108]
unique_ids = set(ids)

print("Unique IDs:", unique_ids)  # {101, 102, 105, 108}
print("Is 105 present?", 105 in unique_ids)  # True
\end{lstlisting}

\textbf{4. Set Algebra}
\begin{lstlisting}
A = {"Alice", "Bob", "Charlie"}
B = {"Bob", "David", "Alice"}

intersection = A & B  # {"Alice", "Bob"}
union = A | B         # {"Alice", "Bob", "Charlie", "David"}

print("Both Cohorts:", intersection)
print("All Unique Cohort Members:", union)
\end{lstlisting}

\textbf{5. Data Science Application: Categorical Encoding Map}
\begin{lstlisting}
labels = ["cat", "dog", "cat", "bird", "dog"]

label2idx = {}
for label in labels:
    if label not in label2idx:
        label2idx[label] = len(label2idx)

print("Label Encoding Map:", label2idx)
# Output: {'cat': 0, 'dog': 1, 'bird': 2}
\end{lstlisting}

\section*{Week 6: Functions, Arguments \& Scope}

\subsection*{Detailed Solutions}

\textbf{1. Basic Function}
\begin{lstlisting}
def is_even(n):
    return n % 2 == 0

print(is_even(4))  # True
print(is_even(7))  # False
\end{lstlisting}

\textbf{2. Default Arguments}
\begin{lstlisting}
def compute_power(base, exponent=2):
    return base ** exponent

print(compute_power(5))     # Output: 25 (5^2)
print(compute_power(2, 3))  # Output: 8  (2^3)
\end{lstlisting}

\textbf{3. Flexible Aggregation (`*args`)}
\begin{lstlisting}
def calculate_mean(*args):
    if not args:
        return 0.0
    return sum(args) / len(args)

print(calculate_mean(10, 20, 30, 40))  # Output: 25.0
print(calculate_mean())                # Output: 0.0
\end{lstlisting}

\textbf{4. Variable Scope}
\begin{lstlisting}
# Output:
# Inside: 100
# Outside: 50
\end{lstlisting}
\textbf{Explanation:} Inside `modify_val()`, the assignment `x = 100` creates a new \emph{local} variable `x` within the function scope. It shadows the global `x` without modifying it. Once the function finishes, the local variable is destroyed, and the global variable `x` remains `50`.

\textbf{5. Data Science Application: Min-Max Scaler Function}
\begin{lstlisting}
def min_max_scaler(data):
    if not data:
        return []
    x_min = min(data)
    x_max = max(data)
    if x_min == x_max:
        return [0.0] * len(data)
    
    return [round((x - x_min) / (x_max - x_min), 3) for x in data]

sample_data = [10, 25, 50, 75, 100]
print(min_max_scaler(sample_data))
# Output: [0.0, 0.167, 0.444, 0.722, 1.0]
\end{lstlisting}

\section*{Week 7: Functional Programming, Comprehensions \& Generators}

\subsection*{Detailed Solutions}

\textbf{1. List Comprehension Filtering}
\begin{lstlisting}
words = ["python", "data", "science", "ai", "machine", "ml"]
filtered = [w.upper() for w in words if len(w) > 3]

print(filtered)  # Output: ['PYTHON', 'DATA', 'SCIENCE', 'MACHINE']
\end{lstlisting}

\textbf{2. Lambda Sorting}
\begin{lstlisting}
students = [("Alice", 88), ("Bob", 95), ("Charlie", 78), ("David", 90)]

sorted_students = sorted(students, key=lambda student: student[1], reverse=True)
print(sorted_students)
# Output: [('Bob', 95), ('David', 90), ('Alice', 88), ('Charlie', 78)]
\end{lstlisting}

\textbf{3. Zipping Vectors}
\begin{lstlisting}
x = [1, 2, 3]
y = [10, 20, 30]

elementwise_sum = [a + b for a, b in zip(x, y)]
print(elementwise_sum)  # Output: [11, 22, 33]
\end{lstlisting}

\textbf{4. Squares Generator}
\begin{lstlisting}
def generate_squares(n):
    for i in range(1, n + 1):
        yield i ** 2

for val in generate_squares(5):
    print(val, end=" ")  # Output: 1 4 9 16 25
print()
\end{lstlisting}

\textbf{5. Data Science Application: Text Preprocessing Pipeline}
\begin{lstlisting}
text = "   Data Science and Machine Learning 101   "

# 1 & 2. Strip whitespace, lower case, tokenize
tokens = text.strip().lower().split()

# 3. Filter out numeric tokens
clean_tokens = list(filter(lambda w: not w.isdigit(), tokens))

print(clean_tokens)
# Output: ['data', 'science', 'and', 'machine', 'learning']
\end{lstlisting}

\section*{Week 8: File Handling \& Exception Management}

\subsection*{Detailed Solutions}

\textbf{1. File Writing \& Reading}
\begin{lstlisting}
# Part A: Writing log lines
with open("log.txt", "w") as f:
    f.write("2026-07-26 INFO Model training started\n")
    f.write("2026-07-26 INFO Epoch 1 complete\n")
    f.write("2026-07-26 INFO Model saved\n")

# Part B: Counting lines
with open("log.txt", "r") as f:
    line_count = sum(1 for _ in f)

print(f"Total lines in log.txt: {line_count}")  # Output: 3
\end{lstlisting}

\textbf{2. Handling Missing Files}
\begin{lstlisting}
try:
    with open("missing_dataset.csv", "r") as f:
        content = f.read()
except FileNotFoundError:
    print("Dataset file not found. Please check filepath.")
\end{lstlisting}

\textbf{3. Safe Type Conversion Function}
\begin{lstlisting}
def safe_float(value):
    try:
        return float(value)
    except (ValueError, TypeError):
        return None

print(safe_float("3.14"))  # 3.14
print(safe_float("abc"))   # None
print(safe_float([1, 2]))  # None
\end{lstlisting}

\textbf{4. Parsing Log File Data}
\begin{lstlisting}
total_score = 0.0
valid_count = 0

with open("scores.txt", "r") as f:
    for line in f:
        try:
            val = float(line.strip())
            total_score += val
            valid_count += 1
        except ValueError:
            # Skip invalid header or corrupted text lines
            continue

if valid_count > 0:
    avg = total_score / valid_count
    print(f"Average Score: {avg:.2f}")
\end{lstlisting}

\textbf{5. Data Science Application: CSV Line Cleaner}
\begin{lstlisting}
def clean_csv(input_path, output_path):
    cleaned_lines = 0
    with open(input_path, "r") as fin, open(output_path, "w") as fout:
        for line in fin:
            stripped = line.strip()
            # Ignore empty lines and comment lines
            if stripped and not stripped.startswith("#"):
                fout.write(line)
                cleaned_lines += 1
    print(f"Cleaned CSV written with {cleaned_lines} valid rows.")
\end{lstlisting}

\section*{Week 9: Recursion \& Searching Algorithms}

\subsection*{Detailed Solutions}

\textbf{1. Recursive Sum}
\begin{lstlisting}
def recursive_sum(n):
    if n <= 1:
        return n
    return n + recursive_sum(n - 1)

print(recursive_sum(10))  # Output: 55
\end{lstlisting}

\textbf{2. Recursive String Reversal}
\begin{lstlisting}
def reverse_string(s):
    if len(s) <= 1:
        return s
    return s[-1] + reverse_string(s[:-1])

print(reverse_string("python"))  # Output: "nohtyp"
\end{lstlisting}

\textbf{3. Binary Search Trace}
Array: `[2, 5, 8, 12, 16, 23, 38, 56, 72, 91]`, Target: `23`.
\begin{enumerate}
    \item \textbf{Step 1:} `low = 0`, `high = 9`. `mid = (0 + 9) // 2 = 4`. `arr[4] = 16`. Since $23 > 16$, search right half: `low = 5`.
    \item \textbf{Step 2:} `low = 5`, `high = 9`. `mid = (5 + 9) // 2 = 7`. `arr[7] = 56`. Since $23 < 56$, search left half: `high = 6`.
    \item \textbf{Step 3:} `low = 5`, `high = 6`. `mid = (5 + 6) // 2 = 5`. `arr[5] = 23`. Match found!
\end{enumerate}
\textbf{Total Comparisons:} $\mathbf{3}$ comparison steps required.

\textbf{4. Recursive Fibonacci with Memoization}
\begin{lstlisting}
def fib_memo(n, memo=None):
    if memo is None:
        memo = {}
    if n in memo:
        return memo[n]
    if n <= 1:
        return n
    
    memo[n] = fib_memo(n - 1, memo) + fib_memo(n - 2, memo)
    return memo[n]

print(fib_memo(50))  # Output: 12586269025 (computes instantly!)
\end{lstlisting}

\textbf{5. Data Science Application: Square Root via Binary Search}
\begin{lstlisting}
def sqrt_binary_search(val, tol=1e-5):
    if val < 0:
        return None
    low, high = 0.0, max(1.0, val)
    
    while high - low > tol:
        mid = (low + high) / 2.0
        if mid * mid < val:
            low = mid
        else:
            high = mid
            
    return round((low + high) / 2.0, 5)

print("sqrt(25) =", sqrt_binary_search(25))  # Output: 5.0
print("sqrt(2) =", sqrt_binary_search(2))    # Output: 1.41421
\end{lstlisting}

\section*{Week 10: Object-Oriented Programming (OOP)}

\subsection*{Detailed Solutions}

\textbf{1. Basic Class Creation}
\begin{lstlisting}
import math

class Circle:
    def __init__(self, radius):
        self.radius = radius

    def area(self):
        return round(math.pi * (self.radius ** 2), 2)

    def circumference(self):
        return round(2 * math.pi * self.radius, 2)

c = Circle(5.0)
print("Area:", c.area())                  # Output: 78.54
print("Circumference:", c.circumference())  # Output: 31.42
\end{lstlisting}

\textbf{2. Class Encapsulation}
\begin{lstlisting}
class BankAccount:
    def __init__(self, balance=0.0):
        self.__balance = balance  # Private attribute

    def deposit(self, amount):
        if amount > 0:
            self.__balance += amount

    def withdraw(self, amount):
        if 0 < amount <= self.__balance:
            self.__balance -= amount
            return True
        return False

    def get_balance(self):
        return self.__balance

acc = BankAccount(100.0)
acc.deposit(50.0)
acc.withdraw(30.0)
print("Balance:", acc.get_balance())  # Output: 120.0
\end{lstlisting}

\textbf{3. Class vs Instance Variables}
\textbf{Class variables} are shared by all instances of a class. \textbf{Instance variables} are bound to a specific object instance.
\begin{lstlisting}
class Dog:
    species = "Canis lupus"  # Class Variable shared across all dogs

    def __init__(self, name):
        self.name = name     # Instance Variable unique to each dog

d1, d2 = Dog("Buddy"), Dog("Max")
print(d1.species, d1.name)  # Canis lupus Buddy
print(d2.species, d2.name)  # Canis lupus Max
\end{lstlisting}

\textbf{4. Inheritance Hierarchy}
\begin{lstlisting}
class Vehicle:
    def __init__(self, brand):
        self.brand = brand

    def drive(self):
        print(f"Driving {self.brand} vehicle...")

class ElectricCar(Vehicle):
    def __init__(self, brand, battery_capacity):
        super().__init__(brand)
        self.battery_capacity = battery_capacity

    def drive(self):  # Method Overriding
        print(f"Driving {self.brand} silently on {self.battery_capacity} kWh battery")

tesla = ElectricCar("Tesla", 85)
tesla.drive()  # Output: Driving Tesla silently on 85 kWh battery
\end{lstlisting}

\textbf{5. Data Science Application: Custom StandardScaler Class}
\begin{lstlisting}
import math

class StandardScaler:
    def __init__(self):
        self.mean = 0.0
        self.std = 1.0

    def fit(self, data):
        self.mean = sum(data) / len(data)
        variance = sum((x - self.mean) ** 2 for x in data) / len(data)
        self.std = math.sqrt(variance)

    def transform(self, data):
        if self.std == 0:
            return [0.0] * len(data)
        return [round((x - self.mean) / self.std, 3) for x in data]

scaler = StandardScaler()
train_data = [10.0, 20.0, 30.0, 40.0, 50.0]
scaler.fit(train_data)

scaled = scaler.transform(train_data)
print("Standardized Data:", scaled)
# Output: [-1.414, -0.707, 0.0, 0.707, 1.414]
\end{lstlisting}

\section*{Week 11: Scientific Computing: NumPy \& Pandas}

\subsection*{Detailed Solutions}

\textbf{1. NumPy Vectorization}
\begin{lstlisting}
import numpy as np

arr = np.arange(1, 11)

# Part A: Boolean mask for odds
mask = (arr % 2 != 0)

# Part B: Extract odds
odds = arr[mask]
print("Odds:", odds)  # [1, 3, 5, 7, 9]

# Part C: Replace evens with -1
arr[arr % 2 == 0] = -1
print("Modified Array:", arr)  # [ 1 -1  3 -1  5 -1  7 -1  9 -1]
\end{lstlisting}

\textbf{2. Matrix Statistics with NumPy}
\begin{lstlisting}
mat = np.arange(1, 10).reshape(3, 3)

total_sum = mat.sum()
col_means = mat.mean(axis=0)
row_maxs = mat.max(axis=1)

print("Total Sum:", total_sum)      # 45
print("Column Means:", col_means)  # [4. 5. 6.]
print("Row Maxima:", row_maxs)     # [3 6 9]
\end{lstlisting}

\textbf{3. Pandas Filtering \& Selection}
\begin{lstlisting}
import pandas as pd

df = pd.DataFrame({
    'City': ['Mumbai', 'Delhi', 'Bangalore', 'Chennai'],
    'Temp': [32, 40, 24, 34],
    'Humidity': [80, 45, 65, 75]
})

filtered_df = df[(df['Temp'] > 30) & (df['Humidity'] > 70)]
print(filtered_df)
# Output:
#       City  Temp  Humidity
# 0   Mumbai    32        80
# 3  Chennai    34        75
\end{lstlisting}

\textbf{4. Pandas Missing Data Handling}
\begin{lstlisting}
import pandas as pd
import numpy as np

df = pd.DataFrame({
    'A': [1.0, 2.0, np.nan, 4.0],
    'B': [np.nan, 10.0, 20.0, 30.0]
})

df['A'] = df['A'].fillna(df['A'].mean())
df['B'] = df['B'].fillna(0.0)

print(df)
# Output:
#           A     B
# 0  1.000000   0.0
# 1  2.000000  10.0
# 2  2.333333  20.0
# 3  4.000000  30.0
\end{lstlisting}

\textbf{5. Data Science Application: Groupby Aggregation}
\begin{lstlisting}
sales_df = pd.DataFrame({
    'Category': ['Tech', 'Tech', 'Office', 'Office', 'Tech'],
    'Revenue': [1200, 800, 150, 300, 2000]
})

total_rev = sales_df.groupby('Category')['Revenue'].sum()
print(total_rev)
# Output:
# Category
# Office     450
# Tech      4000
# Name: Revenue, dtype: int64
\end{lstlisting}

\section*{Week 12: Data Visualization \& Course Summary}

\subsection*{Detailed Solutions}

\textbf{1. Basic Line Plot}
\begin{lstlisting}
import matplotlib.pyplot as plt
import numpy as np

x = np.linspace(-5, 5, 100)
y = x ** 2

plt.figure(figsize=(6, 4))
plt.plot(x, y, color='blue', linewidth=2)
plt.title("Quadratic Function")
plt.xlabel("x")
plt.ylabel("f(x) = x^2")
plt.grid(True)
plt.show()
\end{lstlisting}

\textbf{2. Comparing Distributions (Subplots)}
\begin{lstlisting}
import matplotlib.pyplot as plt
import numpy as np

fig, axes = plt.subplots(1, 2, figsize=(10, 4))

# Subplot 1: Bar Chart
langs = ['Python', 'SQL', 'R']
scores = [95, 80, 60]
axes[0].bar(langs, scores, color='skyblue', edgecolor='black')
axes[0].set_title("Language Popularity")

# Subplot 2: Histogram
data = np.random.uniform(0, 100, 500)
axes[1].hist(data, bins=20, color='orange', edgecolor='black')
axes[1].set_title("Uniform Distribution")

plt.tight_layout()
plt.show()
\end{lstlisting}

\textbf{3. Customizing Scatter Plots}
\begin{lstlisting}
import matplotlib.pyplot as plt
import numpy as np

np.random.seed(42)
x = np.random.rand(20)
y = np.random.rand(20)

plt.scatter(x, y, color='green', marker='^', s=50)
plt.title("Green Triangle Scatter Plot")
plt.xlabel("X values")
plt.ylabel("Y values")
plt.show()
\end{lstlisting}

\textbf{4. Comprehensive Review Question: Data Pipeline Function}
\begin{lstlisting}
import numpy as np

def analyze_dataset(data):
    # 1. Clean negative numbers using list comprehension
    clean_data = [x for x in data if x >= 0]
    
    if not clean_data:
        return {"count": 0, "mean": 0.0, "std": 0.0}
        
    # 2. Convert to NumPy for stats
    arr = np.array(clean_data)
    
    # 3. Return summary dict
    return {
        "count": len(arr),
        "mean": round(float(arr.mean()), 2),
        "std": round(float(arr.std()), 2)
    }

raw_measurements = [12.5, -4.0, 18.2, 22.0, -1.0, 15.8]
print(analyze_dataset(raw_measurements))
# Output: {'count': 4, 'mean': 17.12, 'std': 3.42}
\end{lstlisting}

\textbf{5. Data Science Application: Plotting a Loss Curve}
\begin{lstlisting}
import matplotlib.pyplot as plt

epochs = list(range(1, 11))
loss = [2.5, 1.8, 1.2, 0.8, 0.5, 0.35, 0.25, 0.18, 0.15, 0.12]

plt.figure(figsize=(7, 4))
plt.plot(epochs, loss, 'r-s', label="Training Loss", linewidth=2)
plt.title("Deep Learning Model Loss Curve")
plt.xlabel("Epochs")
plt.ylabel("Cross-Entropy Loss")
plt.legend()
plt.grid(True)
plt.show()
\end{lstlisting}











